\chapter{Proposed Methodology and System Design}
\label{chap:methodology}

\section{Overall Architecture}
The CycleSafe platform is structured as a dual-layer cyber-physical architecture comprising two complementary operational subsystems:
\begin{enumerate}
    \item \textbf{Embedded Physical Telemetry Node (Edge Layer):} An instrumented on-bicycle IoT sensing node that captures real-world lateral ultrasonic clearances, geographic coordinates, 6-axis inertial dynamics, and manual cyclist event confirmations.
    \item \textbf{Microscopic Simulation \& ML Analytics Engine (Cloud/Server Layer):} A containerized backend service running Eclipse SUMO, TraCI, PostgreSQL, Redis, and a supervised XGBoost classification model, exposed via a full-stack web application.
\end{enumerate}

Figure \ref{fig:arch_diagram} visualizes the cross-layer system dataflow and component interconnections.

\begin{figure}[h!]
	\centering
	\setlength{\tabcolsep}{5pt}
	\renewcommand{\arraystretch}{1.2}
	\begin{tabularx}{\linewidth}{|>{\raggedright\arraybackslash}X|>{\raggedright\arraybackslash}X|>{\raggedright\arraybackslash}X|}
		\hline
		\multicolumn{3}{|>{\centering\arraybackslash}p{\dimexpr\linewidth-2\tabcolsep-2\arrayrulewidth}|}{\textbf{LAYER 1: EMBEDDED PHYSICAL TELEMETRY NODE (Edge Hardware)}} \\
		\hline
		\textbf{Ranging \& Kinematics} & \textbf{Microcontroller \& Compute} & \textbf{Power \& User Interface} \\
		Dual JSN-SR04T (20\,Hz) & ESP32 Dual-Core (240\,MHz) & 18650 Li-ion 2500\,mAh (3.7\,V) \\
		MPU-6050 6-Axis IMU & FreeRTOS Dual-Task System & TP4056 Type-C Charger \\
		u-blox NEO-6M GPS & FAT32 MicroSD SPI Logger & AMS1117-3.3V LDO Rail \\
		Handlebar Button (12\,mm) & I2C Bus (GPIO 21/22) & 0.96" SSD1306 OLED \\
		\hline
		\multicolumn{3}{|>{\centering\arraybackslash}p{\dimexpr\linewidth-2\tabcolsep-2\arrayrulewidth}|}{$\Downarrow$ \textit{Standardized Telemetry CSV Schema} $\{t,\, \text{lat},\, \text{lon},\, d_L,\, d_R,\, a_x,\, a_y,\, a_z,\, \text{btn}\} \Downarrow$} \\
		\hline
		\multicolumn{3}{|>{\centering\arraybackslash}p{\dimexpr\linewidth-2\tabcolsep-2\arrayrulewidth}|}{\textbf{LAYER 2: MICROSCOPIC SIMULATION \& ML PLATFORM (Server)}} \\
		\hline
		\textbf{Simulation Subsystem} & \textbf{Machine Learning Engine} & \textbf{Web Application \& Storage} \\
		Eclipse SUMO 1.27.1 + TraCI & XGBoost 3-Class Booster & Go 1.23 REST API Service \\
		10\,km Berlin Arterial Corridor & Trained on SimRA Dataset & PostgreSQL 16 Relational DB \\
		1,244 Stochastic Vehicles & StVO \S\,5(4) Rule Encoders & Redis 7 In-Memory Cache \\
		Sublane Lateral Dynamics (0.2\,m) & Modal Severity Matrix & Leaflet + Chart.js Interface \\
		\hline
	\end{tabularx}
	\caption{End-to-end architectural dataflow of the CycleSafe platform.}
	\label{fig:arch_diagram}
\end{figure}

\section{Design Rationale}
The selection of hardware components and software libraries is guided by strict engineering tradeoffs between measurement accuracy, power consumption, cost, and open-source accessibility:

\begin{itemize}
    \item \textbf{Espressif ESP32 Microcontroller:} Chosen for its dual-core 32-bit architecture operating at 240\,MHz with 520\,KB SRAM. Core 0 is dedicated to high-frequency FreeRTOS sensor sampling (ultrasonic pulse capture and IMU polling), while Core 1 handles NMEA sentence parsing, OLED display rendering, and SPI MicroSD block writes without timing jitter.
    \item \textbf{JSN-SR04T Waterproof Transducers:} Unlike exposed open-mesh ultrasonic sensors (such as HC-SR04), the JSN-SR04T features an enclosed, sealed aluminum transducer probe connected via a 2.5\,m coaxial cable. This design provides IP67 environmental protection against rain, road splashes, and dust, essential for outdoor cycling.
    \item \textbf{Shared I2C Bus Multiplexing:} Both the 0.96" SSD1306 OLED display and the MPU-6050 IMU communicate over the identical physical two-wire interface (\texttt{GPIO 21 = SDA}, \texttt{GPIO 22 = SCL}). Because the OLED controller resides at 7-bit address \texttt{0x3C} and the MPU-6050 resides at \texttt{0x68}, the bus supports concurrent operation with zero additional GPIO pin overhead.
    \item \textbf{Dedicated AMS1117-3.3V Power Regulation Rail:} During write operations, MicroSD cards draw transient current spikes exceeding 100\,mA, while the ESP32 RF subsystem draws up to 240\,mA during Wi-Fi beacon transmission. Supplying these peripherals from a dedicated low-dropout regulator isolated from the raw battery prevents voltage sagging below the 2.7\,V brownout threshold.
\end{itemize}

\section{Algorithms and Operating Principles}

\subsection{Acoustic Time-of-Flight Distance Measurement}
The distance $d_{\text{measured}}$ from the ultrasonic transducer faceplate to the nearest reflecting vehicle surface is governed by the two-way acoustic Time-of-Flight principle:
\begin{equation}
    d_{\text{measured}} = \frac{c \cdot \Delta t}{2}
\end{equation}
where $\Delta t$ is the measured echo pulse duration (round-trip transit time in microseconds), and $c$ is the speed of sound in air. At standard ambient room temperature ($20^\circ\text{C}$), $c \approx 343\,\text{m/s}$, yielding a spatial resolution of approximately $0.17\,\text{mm}$ per microsecond of echo pulse duration.

\subsection{Dynamic Roll-Angle Tilt Compensation}
\label{sec:tilt_comp}
During cornering and pedaling, the bicycle frame tilts at an instantaneous roll angle $\phi$ relative to the vertical. An uncompensated transducer rigidly mounted to the frame projects an inclined acoustic beam toward the ground, recording a slanted hypotenuse distance $d_{\text{measured}}$ rather than the true horizontal ground clearance $d_{\text{true}}$ (Figure \ref{fig:tilt_geometry}).

\begin{figure}[h!]
	\centering
	\fbox{
		\parbox{0.88\linewidth}{
			\vspace{4pt}
			\centering\textbf{\large Dynamic Roll-Angle Tilt Compensation}\\[8pt]
			\begin{itemize}
				\item \textbf{Tilt Angle Measurement ($\phi$):} Monitored continuously via the onboard MPU-6050 3-axis accelerometer.
				\item \textbf{Trigonometric Correction:} The true horizontal clearance to the passing vehicle is calculated as:
				\begin{equation*}
					d_{\text{true}} = d_{\text{measured}} \cdot \cos(\phi)
				\end{equation*}
				\item \textbf{Practical Example:} At a $20^\circ$ lean angle, an uncompensated reading of $1.60\,\text{m}$ corresponds to a true horizontal separation of $1.60 \times \cos(20^\circ) = 1.50\,\text{m}$.
			\end{itemize}
			\vspace{2pt}
		}
	}
	\caption{Geometric principle of dynamic roll-angle tilt compensation.}
	\label{fig:tilt_geometry}
\end{figure}

The true horizontal clearance to the overtaking vehicle is thus determined by:
\begin{equation}
    d_{\text{true}} = d_{\text{measured}} \cdot \cos(\phi)
\end{equation}
where $\phi$ is the instantaneous frame roll angle sampled from the IMU.

\subsection{Machine Learning Conflict Classification Engine}
\label{sec:xgboost_formulation}
Rather than relying solely on a rigid distance threshold of 1.5\,m across all road conditions, CycleSafe utilizes an Extreme Gradient Boosting (\textbf{XGBoost}) multi-class classifier \cite{chen2016xgboost}. 

For each detected encounter, the model extracts a compact feature vector of six observable metrics:
\begin{equation}
    \mathbf{x} = \left[ d_{\text{min}},\, v_{\text{ego}},\, v_{\text{other}},\, a_{\text{ego,min}},\, \Delta v,\, C_{\text{type}} \right]
\end{equation}
where:
\begin{itemize}
    \item $d_{\text{min}}$ is the minimum roll-compensated clearance observed during the encounter [m],
    \item $v_{\text{ego}}, v_{\text{other}}$ are the forward velocities of the bicycle and passing vehicle [m/s],
    \item $a_{\text{ego,min}}$ is the peak braking deceleration recorded by the cyclist [$\text{m/s}^2$],
    \item $\Delta v = v_{\text{other}} - v_{\text{ego}}$ is the relative overtaking velocity [m/s],
    \item $C_{\text{type}}$ is the counterpart vehicle category (bicycle, passenger car, truck/bus, or motorcycle).
\end{itemize}

The XGBoost model constructs an ensemble of shallow decision trees. During training, the algorithm learns optimal split thresholds across the six features to distinguish between safe passes, close passes, and critical incidents. By combining the predictions of multiple decision trees, the model achieves high classification accuracy while resisting noise from momentary sensor dropouts.

\subsection{Statutory Clearances and Vehicle Size Considerations}
The classifier reflects established cycling safety guidelines (such as German StVO \S\,5(4) \cite{stvo2020germany}) while accounting for vehicle size differences:
\begin{itemize}
    \item \textbf{Heavy Vehicles (Trucks and Buses):} Because large commercial vehicles displace large aerodynamic volumes and have significant blind zones, passes below 2.0\,m are categorized with higher hazard weight.
    \item \textbf{Passenger Cars:} Standard 1.5\,m clearance is enforced as the baseline safe passing threshold.
    \item \textbf{Peer Bicycles:} On segregated bicycle paths, cyclists routinely ride side-by-side or pass with smaller clearances (0.8\,m to 1.2\,m) without imminent danger. The model recognizes peer bicycles to avoid false alarms.
\end{itemize}

\section{Data and Empirical Sources}
The models and simulations in CycleSafe are calibrated using two empirical open-access transportation datasets:
\begin{enumerate}
    \item \textbf{TU Berlin SimRA Open Dataset:} Comprises over 25,000 real-world bicycle rides collected across Berlin \cite{simra2020crowdsourcing}. Ground-truth near-miss annotations and accelerometer traces were extracted to train and validate the XGBoost classifier.
    \item \textbf{Berlin Senate Traffic Census 2024/2025 (Oberbaumbr{\"u}cke Counting Station):} Official 12-hour automated traffic counts (07:00--19:00) provided by the \textit{Senatsverwaltung f{\"u}r Mobilit{\"a}t, Verkehr, Klimaschutz und Umwelt} (Document S19-23905 and \textit{Bericht Radverkehr 2025}) \cite{berlin2025report}. The dataset records:
    \begin{itemize}
        \item Passenger Cars and Light Commercial Vehicles (Pkw + Lfw): 12,260 (90.6\% of motor vehicles)
        \item Heavy-Duty Trucks and Buses (Lkw + Bus): 676 (5.0\%)
        \item Motorbikes and Scooters (Krad): 595 (4.4\%)
        \item Commuter Bicycles: 1,397
    \end{itemize}
\end{enumerate}

\section{Ethics, Privacy, and Data Governance}
Because cycling telemetry captures fine-grained spatiotemporal trajectories that could reveal residential addresses or personal daily routines, CycleSafe incorporates privacy-by-design principles aligned with EU General Data Protection Regulation (GDPR) standards:
\begin{itemize}
    \item \textbf{No Optical Surveillance:} The telemetry unit relies entirely on ultrasonic time-of-flight acoustic pulses and inertial sensors; no video cameras or facial recognition systems are deployed.
    \item \textbf{Spatial Obfuscation at Trip Terminals:} In production ingestion, the initial and concluding 200 meters of recorded GPS tracks are clipped to prevent home or workplace deanonymization.
\end{itemize}
